\section{Tests et débogage}

\subsection{Stratégie de tests}

La stratégie accorde la priorité aux invariants backend, car le frontend ne doit
pas être la seule barrière contre une action interdite. La campagne principale
utilise SQLite en mémoire, une queue synchrone, un mailer en mémoire et un
diffuseur nul afin de rester rapide et reproductible. Un second ensemble est
exécuté avec PostgreSQL et Redis dans GitHub Actions afin de vérifier les
migrations, les requêtes spécifiques au moteur et l'infrastructure réellement
utilisée en production. Aucun email réel ni événement WebSocket public n'est
émis pendant ces campagnes.

\begin{longtable}{|p{2.8cm}|p{3.9cm}|p{3.9cm}|p{2.8cm}|}
  \caption{Niveaux de validation appliqués}
  \label{tab:implementation-test-strategy}\\
  \hline
  \textbf{Niveau} & \textbf{Objet} & \textbf{Exemples} &
  \textbf{Preuve}\\
  \hline
  \endfirsthead
  \hline
  \textbf{Niveau} & \textbf{Objet} & \textbf{Exemples} &
  \textbf{Preuve}\\
  \hline
  \endhead
  Tests métier & Invariants, projections et transitions &
  disponibilité, transactions OCPP, tarifs, cycles commerciaux &
  Assertions PHPUnit.\\\hline
  Tests API & Validation, autorisation, sérialisation et effets &
  authentification, stations, alertes, rapports, paiement &
  Réponses et base vérifiées.\\\hline
  Tests d'isolation & Accès croisé entre organisations &
  utilisateurs, actifs, interventions, exports &
  Réponses 403/404 attendues.\\\hline
  Tests d'adaptateur & Contrat avec un service remplaçable &
  WireMock, webhook signé, Google OAuth simulé &
  Doubles et serveur local.\\\hline
  Tests OCPP & Lecture des profils et scénarios du simulateur &
  connexion, plug/unplug, transactions, projections &
  Node Test et PHPUnit.\\\hline
  Qualité frontend & Types, imports et règles statiques &
  compilation Vite et Oxlint &
  Commandes sans erreur.\\\hline
  Vérification visuelle & Mise en page et parcours représentatifs &
  dashboards, stations, cartes, formulaires, documents &
  Inspection navigateur/PDF.\\
  \hline
\end{longtable}

\subsection{Résultats mesurés}

La campagne d'intégration continue exécutée le 13 août 2026 produit les résultats
suivants :

\begin{table}[H]
  \centering
  \caption{Résultats de validation du dépôt}
  \label{tab:implementation-test-results}
  \begin{tabularx}{\textwidth}{|p{4.5cm}|p{3.2cm}|Y|}
    \hline
    \rowcolor{ctmint}
    \textbf{Contrôle} & \textbf{Résultat} & \textbf{Portée}\\\hline

    Tests backend complets &
    265 tests, 2\,761 assertions & Tests Feature et unitaires chargés par
    PHPUnit avec SQLite.\\\hline
    Compatibilité PostgreSQL &
    20 tests, 153 assertions & Migrations, requêtes et contraintes exécutées
    avec PostgreSQL~18.\\\hline
    \texttt{python -m pytest} dans \texttt{ocpp-gateway} &
    23 tests, 0 échec & Messages entrants, commandes distantes, signatures et
    identifiants d'événements.\\\hline
    Tests frontend &
    39 tests, 0 échec & Composants, contrats, utilitaires et comportements
    critiques, complétés par Oxlint et la compilation Vite.\\\hline
    Tests infrastructure et outils &
    39 tests, 0 échec & Compose, CI/CD, sauvegarde, politiques de secrets,
    Simulation Lab et outils OCPP.\\\hline
    Construction Docker &
    Succès & Images de l'API, du frontend, de la passerelle et des simulateurs.\\\hline

  \end{tabularx}
\end{table}

La réussite de la compilation ne prouve pas seule la qualité UX, et le nombre de
tests ne remplace pas une recette. Elle montre néanmoins que les contrats
principaux, les migrations et la majorité des cycles de vie restent cohérents
après l'intégration des développements successifs.

\subsection{Anomalies représentatives et non-régression}

Plusieurs anomalies rencontrées pendant la réalisation ont renforcé les contrôles :

\begin{itemize}[leftmargin=7mm]
  \item une requête depuis l'adresse locale du réseau pouvait recevoir une erreur
  419 lorsque les domaines de session et les origines CSRF ne couvraient que
  \texttt{localhost}. La configuration a été rendue cohérente pour l'hôte local et
  l'adresse LAN de développement ;
  \item une double soumission involontaire d'une invitation déclenchait une erreur
  429. Le bouton est verrouillé pendant la requête et les limites restent actives
  côté serveur ;
  \item un tri \texttt{created\_at} combiné à un \texttt{GROUP BY} était accepté
  par certains moteurs mais refusé par PostgreSQL. La requête d'agrégation des
  photos d'intervention a été corrigée et couverte par un test ;
  \item \texttt{crypto.randomUUID()} échouait sur un navigateur mobile servi en
  HTTP non sécurisé. L'identifiant temporaire utilise désormais une solution
  compatible qui ne modifie pas l'identité métier du backend ;
  \item des URL WebSocket figées sur \texttt{localhost} empêchaient Reverb de
  fonctionner depuis un téléphone. Le frontend construit l'hôte de développement
  à partir de sa configuration plutôt que de supposer une machine unique.
\end{itemize}

Les réponses d'erreur de l'API sont normalisées pour ne pas exposer une trace SQL,
un chemin local ou un secret. En environnement de production,
\texttt{APP\_DEBUG=false} reste obligatoire.

\section{Intégration des composants}

\subsection{Intégration bout en bout}

L'intégration verticale est vérifiable sans accès direct du frontend à la borne
ou au fournisseur de paiement. React appelle uniquement l'API publique Laravel.
Laravel crée les intentions métier et les commandes. La passerelle OCPP échange
avec le simulateur puis renvoie des événements signés. PostgreSQL conserve les
preuves et projections. Reverb diffuse les changements utiles à l'interface.
WireMock représente un système financier externe indépendant.

\begin{figure}[H]
  \centering
  \begin{tikzpicture}[
    node distance=7mm and 8mm,
    box/.style={draw=ctline,fill=ctpaper,rounded corners=2mm,
      minimum width=2.7cm,minimum height=1.0cm,align=center,
      font=\small\bfseries,text=ctdark},
    arrow/.style={-{Latex[length=2mm]},thick,draw=ctgreen}
  ]
    \node[box] (react) {React};
    \node[box,right=of react] (laravel) {Laravel};
    \node[box,right=of laravel] (db) {PostgreSQL};
    \node[box,below=of laravel] (gateway) {Passerelle OCPP};
    \node[box,left=of gateway] (simulator) {Simulateur SAP};
    \node[box,right=of gateway] (payment) {WireMock};
    \draw[arrow] (react) -- node[above,font=\scriptsize]{HTTP} (laravel);
    \draw[arrow] (laravel) -- node[above,font=\scriptsize]{transactions} (db);
    \draw[arrow] (laravel) -- node[right,font=\scriptsize]{commandes} (gateway);
    \draw[arrow] (gateway) -- node[below,font=\scriptsize]{WebSocket OCPP} (simulator);
    \draw[arrow] (laravel) -- node[above right,font=\scriptsize]{paiement} (payment);
    \draw[arrow,dashed] (laravel) to[bend left=20]
      node[above,font=\scriptsize]{Reverb} (react);
  \end{tikzpicture}
  \caption{Composants traversés par le MVP vertical}
  \label{fig:implementation-vertical-integration}
\end{figure}

Cette structure facilite le diagnostic. Un écran n'actualisant pas une donnée
peut être analysé successivement au niveau de la requête React, de la réponse API,
de la projection PostgreSQL, de l'événement reçu par la passerelle et du message
émis par le simulateur.

\subsection{Services externes intégrés}

\begin{table}[H]
  \centering
  \caption{Intégrations techniques et responsabilités}
  \label{tab:external-integrations}
  \begin{tabularx}{\textwidth}{|p{3.0cm}|p{3.1cm}|Y|Y|}
    \hline
    \rowcolor{ctmint}
    \textbf{Intégration} & \textbf{Technologie} & \textbf{Responsabilité} &
    \textbf{Frontière}\\\hline
    Identité sociale & Google OAuth 2.0 & Vérifier l'identité et l'email &
    Les rôles et organisations restent locaux.\\\hline
    Email transactionnel & Resend & Envoyer vérifications, invitations et alertes &
    Les secrets et l'envoi restent côté serveur.\\\hline
    Cartographie & Leaflet / OpenStreetMap & Afficher tuiles, marqueurs et itinéraires &
    La position courante du client n'est pas persistée.\\\hline
    Paiement & Adaptateur de paiement & Autoriser, capturer, libérer ou débiter &
    Le fournisseur ne modifie jamais directement une session.\\\hline
    Simulation du paiement & WireMock 3 & Reproduire succès, refus, délai et webhook &
    Aucun moyen de paiement réel n'est utilisé.\\\hline
    Simulation des bornes & Simulateur SAP & Produire des connexions et messages OCPP &
    Le simulateur ne lit ni l'API métier ni la base.\\\hline
  \end{tabularx}
\end{table}

\section{Validation et vérification}

\subsection{Limites de la validation}

\begin{scopebox}[Frontière des résultats]
\begin{itemize}[leftmargin=7mm]
  \item Le simulateur SAP valide les messages OCPP 1.6J, mais pas le comportement
  électrique, le modem ou le firmware d'une borne réelle.
  \item WireMock valide le contrat HTTP et les webhooks, mais aucune transaction
  bancaire ni exigence PCI DSS.
  \item Les tests automatisés n'ont pas encore été complétés par une campagne de
  charge OCPP multi-borne ni par des tests de performance formels.
  \item Le déploiement public, TLS, CI/CD et les sauvegardes sont opérationnels,
  mais ils n'ont pas encore fait l'objet d'une mesure de disponibilité de longue
  durée, d'un test de charge représentatif ni d'un exercice complet de reprise.
  \item La recette fonctionnelle formelle et son procès-verbal restent à réaliser
  avec le représentant habilité de l'entreprise.
  \item OCPP 2.0.1, Microsoft OAuth, le paiement réel et le firmware restent hors
  du MVP validé.
\end{itemize}
\end{scopebox}

\subsection{Évolution par rapport au cahier des charges initial}
\label{subsec:scope-evolution}

Le \emph{Cahier des Charges Fonctionnel -- Développement d'une plateforme Web
de Gestion des Bornes de Recharge Électrique Connectées}, version 1.0, constitue
la référence initiale transmise par l'entreprise \cite{tacticCahierCharges2026}.
Il décrit un produit cible large et conserve volontairement plusieurs
alternatives technologiques. Le Product Backlog et les revues de Sprint ont
transformé cette intention en un MVP vérifiable sans modifier rétroactivement
le document d'origine.

\begin{longtable}{|p{3.0cm}|p{4.5cm}|p{6.8cm}|}
  \caption{Principales adaptations du périmètre initial}
  \label{tab:scope-main-adaptations}\\
  \hline
  \textbf{Décision} & \textbf{Prévision initiale} & \textbf{Choix retenu et justification}\\
  \hline
  \endfirsthead
  \hline
  \textbf{Décision} & \textbf{Prévision initiale} & \textbf{Choix retenu et justification}\\
  \hline
  \endhead
  Socle backend & Laravel 12, PHP 8.3 et JWT &
  Laravel 13 sur PHP 8.5 local, avec compatibilité Composer à partir de PHP 8.3.
  Sanctum protège la SPA par session et CSRF, sans jeton JWT stocké dans le
  navigateur.\\\hline
  Traitements internes & Redis pour cache et files &
  Redis~7 est utilisé pour le cache, les sessions et les files d'attente, tandis
  que PostgreSQL reste réservé aux données métier persistantes.\\\hline
  Interface analytique & React, Ant Design, TypeScript et ECharts &
  React, Ant Design et TypeScript sont conservés ; Recharts remplace ECharts
  dans les tableaux de bord, sans changer les besoins métier.\\\hline
  Protocole de borne & OCPP 1.6 et OCPP 2.0.1 &
  OCPP 1.6J est implémenté et testé avec la passerelle et le simulateur SAP.
  OCPP 2.0.1 est reporté afin de stabiliser d'abord le parcours vertical.\\\hline
  Paiement & Carte, portefeuille, paiement différé et remboursement &
  Un contrat de paiement remplaçable est validé avec WireMock et des méthodes
  carte, e-DINAR et D17 simulées. Aucun débit bancaire ni remboursement réel
  n'est revendiqué.\\\hline
  Équipement physique & Réseau de bornes connectées &
  Le commissioning, l'authentification et les échanges sont démontrés avec un
  simulateur OCPP. La validation électrique, modem et firmware exige une borne
  physique distincte.\\\hline
  Modules différés & Véhicules, RFID complet et firmware &
  L'identifiant OCPP virtuel et le QR d'entrée couvrent le démarrage du MVP.
  Les profils véhicule, les badges RFID physiques et la gestion distante du
  firmware sont exclus de la version validée.\\
  \hline
\end{longtable}

La matrice détaillée de l'annexe~\ref{app:specification-coverage} reprend les
18 modules, les messages OCPP, les choix d'architecture, les exigences non
fonctionnelles, les interfaces et les livrables. Elle emploie cinq statuts :
\emph{réalisée}, \emph{adaptée}, \emph{partielle}, \emph{reportée} et
\emph{à valider en production}. Neuf modules sont réalisés ou adaptés au MVP,
sept sont partiels et deux sont reportés. Cette lecture distingue ainsi un
écart assumé d'une fonctionnalité simplement oubliée.

\subsection{Traçabilité avec le backlog et les Sprints}

\begin{longtable}{|C{1.5cm}|p{3.3cm}|p{4.5cm}|p{4.0cm}|}
  \caption{Preuves de réalisation par Sprint}
  \label{tab:implementation-traceability}\\
  \hline
  \textbf{Sprint} & \textbf{User stories} & \textbf{Preuves dans le code} &
  \textbf{Tests principaux}\\
  \hline
  \endfirsthead
  \hline
  \textbf{Sprint} & \textbf{User stories} & \textbf{Preuves dans le code} &
  \textbf{Tests principaux}\\
  \hline
  \endhead
  S01 & US-27, 39 et travaux de conception &
  Backlog, cas d'utilisation global, architecture et prototype &
  Relectures, validation des parcours et démonstration du prototype.\\\hline
  S02 & US-01--04, 06 &
  Auth, Google, invitations, démo, onboarding, policies &
  Auth, OAuth, onboarding, utilisateurs, organisations.\\\hline
  S03 & US-07--13, 37 &
  Stations, connecteurs, cartes, commissioning, documents &
  Station API, commissioning, documents.\\\hline
  S04 & US-14--16, 19 &
  Gateway, ingestion, projection, commandes, télémétrie &
  OCPP gateway, flotte, supervision, disponibilité, télémétrie.\\\hline
  S05 & US-24--26, 28--30 &
  Tentatives, sessions, transactions, paiements, reçus &
  Recharge distante, session/paiement, webhook WireMock.\\\hline
  S06 & US-20--23, 40 &
  Alertes, interventions, maintenance, SLA, notifications &
  Alertes, rapports terrain, maintenance, notifications.\\\hline
  S07 & US-31, 32, 41, 42 &
  Tarifs, plans, abonnements clients et SaaS &
  Tarifs, plans, commercial organisationnel.\\\hline
  S08 & US-05, 17, 18, 33--36 &
  Dashboards, reporting, audit, paramètres, intégrations &
  Dashboard, rapports, gouvernance, recherche.\\\hline
  S09 & Stabilisation et livrables &
  Exports, documentation, cahier de tests et préparation de la recette &
  Compilation, contrôles transversaux et validation des documents.\\\hline
  S10 & Industrialisation et démonstration publique &
  Docker unifié, Simulation Lab, Redis, CI/CD, Azure, domaine et sauvegardes &
  Tests de politique, construction des images, déploiement et contrôles de santé.\\
  \hline
\end{longtable}

\section{Optimisation des performances}

Les listes volumineuses utilisent pagination et filtres côté serveur. Les
relations nécessaires aux ressources API sont chargées explicitement afin de
limiter les requêtes répétitives. Les agrégats de tableaux de bord sont calculés
par période et périmètre plutôt que reconstruits dans le navigateur. Les travaux
lents, notamment les emails et certaines notifications, passent par la queue.

Pour OCPP, la passerelle conserve la connexion WebSocket mais délègue rapidement
l'événement à Laravel. Le recalcul de disponibilité est idempotent et protégé
contre le chevauchement. Les événements et mesures restent horodatés, ce qui
permet d'appliquer un délai de non-réponse sans dépendre de l'heure du navigateur.

Le frontend applique un découpage par route et charge à la demande les fonctions
coûteuses, notamment la visualisation PDF et le scanner QR. La compilation de
production est contrôlée par la CI. Une campagne Lighthouse sur réseau mobile,
des budgets de bundle et un test de charge multi-stations restent cependant à
formaliser avant de fixer des objectifs de performance contractuels.

\section{Documentation technique}

Le dépôt contient les scripts d'installation, les configurations d'exemple, les
procédures de déploiement et de sauvegarde, les commandes d'exploitation des
simulateurs et le présent rapport. Scramble génère
une spécification OpenAPI 3.1 versionnée dans
\path{docs/api/openapi.json} et une interface d'exploration protégée sous
\path{/docs/api}. Les routes internes OCPP et paiement restent volontairement
exclues de cette documentation publique.

Le manuel utilisateur, le manuel administrateur et le cahier de tests sont
également compilés et versionnés. Le procès-verbal de recette ne sera préparé
qu'au moment de la campagne fonctionnelle avec l'entreprise, puisqu'il devra
reprendre les résultats réellement observés, les réserves et la décision du
responsable habilité.
